\documentclass[a4paper,12pt]{article}

\usepackage{float}
\usepackage{upgreek}
\usepackage{amsmath}
\usepackage{amssymb}
\usepackage{graphicx}
\usepackage{hyperref}
\usepackage[a4paper, margin=2.5cm]{geometry}
\usepackage[utf8]{inputenc}
\usepackage[T1]{fontenc}
\usepackage{lmodern}
\usepackage{xcolor}
\usepackage{listings}

\lstset{
    basicstyle=\ttfamily\small,
    breaklines=true,
    frame=single,
    backgroundcolor=\color{gray!10},
    keywordstyle=\color{blue},
    commentstyle=\color{green!60!black},
}

\renewcommand{\figurename}{Slika}

\begin{document}

    \begin{titlepage}
        \begin{center}
            \vspace*{3cm}
            {\LARGE \textbf{Sferična aberacija sferičnih leč in raytracing}} \\[1cm]
            {\large Seminarska naloga pri predmetu Numerične metode} \\[2cm]

            {\Large Avtor: Taj Šobot} \\[0.5cm]
            {\Large Mentor: doc. dr. Rene Markovič} \\[0.5cm]

            \vfill
            {\large \today}
        \end{center}
    \end{titlepage}

    \flushleft

    % ─────────────────────────────────────────────────────────────────────────
    \section{Uvod}
    % ─────────────────────────────────────────────────────────────────────────

    Sferična aberacija je ena izmed optičnih napak, ki nastane, ko leča ne lomi vseh vpadnih žarkov v isto gorišče.
    V tej nalogi smo analizirali lom žarkov skozi sferično loečo, kjer je ta pojav bistveno opazljiv.
    Bolj oddaljeni so vpadni žarki od optične osi leče, večji pogrešek imajo.
    Rezultat je popačena slika.
    V inštrumentih kot so kamere in teleskopi, kjer nastopajo sistemi leč pa to prevede v zamegljene in ne ostre slike.

    Lečo bomo modelirali kot krogelno ploskev z določenim lomnim količnikom $n$ in polmerom $r$.
    Uporabljali bomo program napisan v C++ programskem jeziku, ki bo z uporabo metode "ray tracing", simuliral prehod svetlobnih žarkov skozi tako lečo.


    % ─────────────────────────────────────────────────────────────────────────
    \section{Fizikalno ozadje}
    % ─────────────────────────────────────────────────────────────────────────

    \subsection{Snellov zakon}

    Ko svetlobni žarek prehaja iz medija z lomnim količnikom $n_1$ v medij z lomnim količnikom
    $n_2$, velja Snellov zakon:
    \begin{equation}
        n_1 \sin\theta_1 = n_2 \sin\theta_2,
        \label{eq:snell}
    \end{equation}
    kjer je $\theta_1$ vpadni kot in $\theta_2$ lomni kot, oba merjena glede na normalo ploskve
    v točki vpada.

    \subsection{Vektorska oblika Snellovega zakona}

    V ray tracingu je priročno Snellov zakon izraziti vektorsko, saj poznamo smer vpadnega žarka
    $\hat{d}$ in normalo ploskve $\hat{n}$, ne pa kotov direktno.

    \subsubsection{Izpeljava}

    Označimo:
    \begin{itemize}
        \item $\hat{d}$ — enotski vektor vpadnega žarka (kaže \emph{v} ploskev),
        \item $\hat{n}$ — enotski normala ploskve (kaže \emph{stran} od ploskve, proti vpadnemu žarku),
        \item $\hat{t}$ — enotski vektor lomljene smeri.
    \end{itemize}

    Kosinuse kotov dobimo iz skalarnih produktov:
    \begin{equation}
        \cos\theta_1 = -\hat{d} \cdot \hat{n}, \qquad \cos\theta_2 = \sqrt{1 - \sin^2\theta_2}.
    \end{equation}

    Iz Snellovega zakona \eqref{eq:snell} sledi:
    \begin{equation}
        \sin^2\theta_2 = \left(\frac{n_1}{n_2}\right)^2 \sin^2\theta_1
        = \left(\frac{n_1}{n_2}\right)^2 (1 - \cos^2\theta_1).
    \end{equation}

    Lom\-ljeni žarek $\hat{t}$ razstavimo na tangencialno in normalno komponento glede na ploskev:
    \begin{align}
        \hat{t} &= \hat{t}_\parallel + \hat{t}_\perp.
    \end{align}

    Tangencialna komponenta je vzporedna z vpadnim žarkom (brez normalne sestavine):
    \begin{equation}
        \hat{t}_\parallel = \frac{n_1}{n_2}\left(\hat{d} + \cos\theta_1\,\hat{n}\right).
    \end{equation}

    Normalna komponenta je vzdolž normale:
    \begin{equation}
        \hat{t}_\perp = -\cos\theta_2\,\hat{n}.
    \end{equation}

    Skupaj torej:
    \begin{equation}
        \boxed{
            \hat{t} = \frac{n_1}{n_2}\,\hat{d}
            + \left(\frac{n_1}{n_2}\cos\theta_1 - \cos\theta_2\right)\hat{n}
        }
        \label{eq:refract_vec}
    \end{equation}
    kjer je $\cos\theta_1 = -\hat{d}\cdot\hat{n}$ in $\cos\theta_2 = \sqrt{1 - (n_1/n_2)^2(1-\cos^2\theta_1)}$.

    \subsubsection{Popolni odboj}

    Kadar je $\sin^2\theta_2 > 1$, lomljenega žarka ni — nastopi popolni notranji odboj (ang.
    \emph{total internal reflection}, TIR). Pogoj:
    \begin{equation}
        \left(\frac{n_1}{n_2}\right)^2(1 - \cos^2\theta_1) > 1.
    \end{equation}
    V tem primeru žarka ne propagiramo naprej.

    % ─────────────────────────────────────────────────────────────────────────
    \section{Geometrija sfere}
    % ─────────────────────────────────────────────────────────────────────────

    \subsection{Enačba sfere}

    Sfera s središčem $\vec{c}$ in polmerom $R$ je množica vseh točk $\vec{P}$, ki zadostijo:
    \begin{equation}
        |\vec{P} - \vec{c}|^2 = R^2.
    \end{equation}

    \subsection{Presečišče žarka s sfero}

    Žarek opišemo parametrično:
    \begin{equation}
        \vec{P}(t) = \vec{o} + t\,\hat{d},
    \end{equation}
    kjer je $\vec{o}$ izhodišče žarka in $\hat{d}$ enotska smer. Vstavimo v enačbo sfere:
    \begin{equation}
        |\vec{o} + t\,\hat{d} - \vec{c}|^2 = R^2.
    \end{equation}

    Označimo $\vec{oc} = \vec{o} - \vec{c}$ in razvijemo:
    \begin{align}
        |\hat{d}|^2\,t^2 + 2(\hat{d}\cdot\vec{oc})\,t + |\vec{oc}|^2 - R^2 &= 0 \\
        t^2 + 2(\hat{d}\cdot\vec{oc})\,t + |\vec{oc}|^2 - R^2 &= 0, \qquad (|\hat{d}|^2 = 1)
    \end{align}

    kar je kvadratna enačba $at^2 + bt + c = 0$ s koeficienti:
    \begin{equation}
        a = 1, \quad b = 2\,\hat{d}\cdot\vec{oc}, \quad c = |\vec{oc}|^2 - R^2.
    \end{equation}

    Diskriminanta:
    \begin{equation}
        D = b^2 - 4c.
    \end{equation}

    \begin{itemize}
        \item Če $D < 0$: žarek ne seka sfere.
        \item Če $D \geq 0$: dve rešitvi $t_{1,2} = \dfrac{-b \pm \sqrt{D}}{2}$.
    \end{itemize}

    Izberemo najmanjši pozitivni $t$ (najbližje presečišče pred žarkom):
    \begin{equation}
        t = \min\{t_1, t_2 \mid t > \varepsilon\},
    \end{equation}
    kjer je $\varepsilon > 0$ majhna vrednost, ki prepreči samopresečišče (ang. \emph{self-intersection}).

    \subsection{Normala v točki presečišča}

    Ko dobimo točko presečišča $\vec{P}_{hit} = \vec{o} + t\,\hat{d}$, je zunanja normala sfere:
    \begin{equation}
        \hat{n} = \frac{\vec{P}_{hit} - \vec{c}}{|\vec{P}_{hit} - \vec{c}|} = \frac{\vec{P}_{hit} - \vec{c}}{R}.
    \end{equation}

    Pri izstopu iz sfere (druga površina) moramo normalo obrniti: $\hat{n} \to -\hat{n}$, da pravilno
    označimo prehod iz gostejšega v redkejši medij.

    % ─────────────────────────────────────────────────────────────────────────
    \section{Sferična aberacija}
    % ─────────────────────────────────────────────────────────────────────────

    \subsection{Paraksijska aproksimacija}

    Za žarke blizu optične osi (majhen $\theta_1$) velja $\sin\theta \approx \theta$, kar dá
    enostavno lensmaker's enačbo:
    \begin{equation}
        \frac{1}{a} + \frac{1}{b} = \frac{1}{f},
    \end{equation}
    kjer je $a$ razdalja predmeta, $b$ razdalja slike in $f$ goriščna razdalja leče.

    \subsection{Izvor aberacije}

    Za žarke daleč od osi paraksijska aproksimacija ne velja več. Dejansko gorišče žarka, ki
    vpade na rob sferne leče z polmerom $R$ in lomnim količnikom $n$, je bliže leči kot paraksijsko
    gorišče. Razlika med paraksijskim in dejanskim gorilnim mestom tvori \emph{sferno aberacijo}.

    Za sferično lečo (obe površini sta krogelni s polmerom $R$) je goriščna razdalja v paraksijski
    aproksimaciji:
    \begin{equation}
        \frac{1}{f} = (n - 1)\left(\frac{1}{R_1} - \frac{1}{R_2}\right),
    \end{equation}
    kar je \emph{lensmaker's enačba} za tanko lečo. Za sferno kroglo z $R_1 = R$ in $R_2 = -R$:
    \begin{equation}
        \frac{1}{f} = \frac{2(n-1)}{R}.
    \end{equation}

    % ─────────────────────────────────────────────────────────────────────────
    \section{Implementacija v GLSL}
    % ─────────────────────────────────────────────────────────────────────────

    \subsection{Compute shader — presečišče s sfero}

    Vsak invokacijski thread obdeluje en piksel izhodne slike. Žarek je metamo vzporedno z
    optično osjo ($z$-smeri), izhodišče žarka pa ustreza položaju piksla v ravnini senzorja.

    \begin{lstlisting}[language=C, caption={Presečišče žarka s sfero v GLSL}]
float intersectSphere(vec3 origin, vec3 dir, vec3 center, float r) {
    vec3 oc = origin - center;
    float b = 2.0 * dot(dir, oc);
    float c = dot(oc, oc) - r * r;
    float D = b * b - 4.0 * c;        // a = 1 ker je |dir| = 1
    if (D < 0.0) return -1.0;
    float t1 = (-b - sqrt(D)) / 2.0;
    float t2 = (-b + sqrt(D)) / 2.0;
    if (t1 > 0.001) return t1;
    if (t2 > 0.001) return t2;
    return -1.0;
}
    \end{lstlisting}

    \subsection{Compute shader — vektorski lom}

    Implementacija enačbe \eqref{eq:refract_vec}:

    \begin{lstlisting}[language=C, caption={Vektorski Snellov zakon v GLSL}]
vec3 refractRay(vec3 d, vec3 normal, float n1, float n2) {
    float cosI  = -dot(normal, d);
    float sinT2 = (n1/n2)*(n1/n2) * (1.0 - cosI*cosI);
    if (sinT2 > 1.0) return vec3(0.0);   // TIR
    float cosT = sqrt(1.0 - sinT2);
    return (n1/n2)*d + (n1/n2*cosI - cosT)*normal;
}
    \end{lstlisting}

    \subsection{Potek žarka skozi lečo}

    Vsak žarek preide skozi dve površini sfere — vstopno in izstopno. Na vsaki površini
    apliciramo lom:

    \begin{lstlisting}[language=C, caption={Zanka dveh prehodov skozi sfero}]
for (int i = 0; i < 2; i++) {
    float t = intersectSphere(rayOrigin, rayDir, lensCenter, u_r);
    if (t <= 0.0) break;

    vec3 hit    = rayOrigin + t * rayDir;
    vec3 normal = normalize(hit - lensCenter);
    if (i == 1) normal = -normal;  // izstop: normala obrnjena navznoter

    float n1 = (i == 0) ? 1.0 : u_n;
    float n2 = (i == 0) ? u_n : 1.0;

    rayDir    = refractRay(rayDir, normal, n1, n2);
    rayOrigin = hit + rayDir * 0.001;  // epsilon odmik
}
    \end{lstlisting}

    Na prvi površini ($i=0$) gre žarek iz zraka ($n_1=1$) v steklo ($n_2=n$), na drugi ($i=1$)
    pa iz stekla ($n_1=n$) nazaj v zrak ($n_2=1$). Normala na izstopu je obrnjena, ker mora
    pri vektorski obliki Snellovega zakona normala kazati \emph{stran od medija, v katerem je
    vpadni žarek}.

    \subsection{Vzorčenje slike na zaslonu}

    Po izstopu iz leče žarek propagiramo do ravnine slike na razdalji $u\_b + u\_a$:

    \begin{lstlisting}[language=C, caption={Vzorčenje objektne ravnine}]
float t_plate  = (objectPos.z - rayOrigin.z) / rayDir.z;
vec3  hit      = rayOrigin + t_plate * rayDir;

// UV koordinate za vzorcenje vhodne teksture
vec2 uv = (hit.xy - objectPos.xy) / u_ObjScale;
if (all(greaterThanEqual(uv, vec2(0.0))) &&
    all(lessThanEqual(uv, vec2(1.0))))
    color = texture(u_input, uv);
else
    color = vec4(0.2 * normalize(hit.xyy), 1.0);
    \end{lstlisting}

    % ─────────────────────────────────────────────────────────────────────────
    \section{Parametri programa}
    % ─────────────────────────────────────────────────────────────────────────

    Program sprejema naslednje uniforme, ki jih med izvajanjem nastavljamo z bližnjicami:

    \begin{table}[H]
        \centering
        \begin{tabular}{|l|l|l|}
            \hline
            \textbf{Uniform} & \textbf{Pomen} & \textbf{Tippična vrednost} \\
            \hline
            \texttt{u\_r}        & Polmer sferne leče              & $0.5 \ldots 2.0$ \\
            \texttt{u\_n}        & Lomni količnik stekla           & $1.5$ (crown glass) \\
            \texttt{u\_a}        & Razdalja predmeta od leče       & $10.0$ (vzporedni žarki $\to \infty$) \\
            \texttt{u\_b}        & Razdalja slike od leče          & nastavljivo \\
            \texttt{u\_xLensPos} & Horizontalni položaj leče       & $0.0$ \\
            \texttt{u\_yLensPos} & Vertikalni položaj leče         & $0.0$ \\
            \texttt{u\_xObjPos}  & Horizontalni položaj objekta    & $0.0$ \\
            \texttt{u\_yObjPos}  & Vertikalni položaj objekta      & $0.0$ \\
            \texttt{u\_ObjScale} & Velikost objekta na zaslobu     & $1.0$ \\
            \hline
        \end{tabular}
        \caption{Uniformni parametri compute shaderja}
    \end{table}

    % ─────────────────────────────────────────────────────────────────────────
    \section{Zaključek}
    % ─────────────────────────────────────────────────────────────────────────

    Implementirali smo ray tracer za simulacijo sferične aberacije, ki deluje v realnem času
    na GPU. Vsak piksel izhodne slike ustreza enemu svetlobnemu žarku, ki ga propagiramo
    skozi sferno lečo z dvema lomoma (vstop in izstop), pri čemer na vsaki površini apliciramo
    vektorsko obliko Snellovega zakona. Sferična aberacija se jasno pokaže kot razmazanost
    slike, ki se povečuje z večjim polmerom žarka in večjim lomnim količnikom. Interaktivna
    nastavitev parametrov omogoča intuitivno razumevanje, kako posamezni parametri vplivajo
    na kakovost slike.

    \pagebreak
    \section{Viri}

    \begin{thebibliography}{9}
        \bibitem{hecht}
        E. Hecht, \textit{Optics}, 5th ed., Pearson, 2017.

        \bibitem{shirley}
        P. Shirley, \textit{Ray Tracing in One Weekend}, 2020. \\
        Dostopno na: \url{https://raytracing.github.io/}

        \bibitem{glsl}
        Khronos Group, \textit{GLSL Specification 4.60}. \\
        Dostopno na: \url{https://registry.khronos.org/OpenGL/specs/gl/GLSLangSpec.4.60.pdf}

        \bibitem{opengl}
        Khronos Group, \textit{OpenGL 4.3 Core Profile Specification}, 2012. \\
        Dostopno na: \url{https://registry.khronos.org/OpenGL/specs/gl/glspec43.core.pdf}
    \end{thebibliography}

\end{document}